\section{Introduction}\label{sec:intro}

\paragraph{Problem.}
Diagnostic models for skin lesions are typically trained on curated dermoscopic images, yet real
consumer-device photographs span a wide quality range --- motion blur, poor illumination, incomplete
framing, colour casts, and low contrast --- on which a model accurate on clean data can be both wrong
and confident. The deployment question is therefore not only ``what is the diagnosis?'' but ``is this
image good enough to diagnose, and if not, what should the system do about it?''

\paragraph{Why existing tools fall short.}
Three families of methods each address part of the problem but none closes the loop.
(i) \emph{Post-hoc calibration} --- temperature scaling \citep{guo2017calibration}, Platt and isotonic
regression --- adjusts confidence after training but is quality-oblivious, unable to separate
uncertainty from genuine ambiguity from uncertainty due to degradation.
(ii) \emph{Uncertainty estimation} via MC dropout \citep{gal2016dropout} or deep ensembles
\citep{lakshminarayanan2017simple} likewise yields quality-oblivious uncertainty at added inference
cost.
(iii) \emph{Generic image enhancement} (e.g.\ Real-ESRGAN \citep{wang2021realesrgan}, Restormer
\citep{zamir2022restormer}) improves perceptual quality with no guarantee that diagnostic content is
preserved, and generative restoration risks inventing clinically unacceptable structures (pigment
networks, vessels).

\paragraph{Closed-loop quality-triage.}
We argue that quality must be a first-class signal threaded through the whole pipeline. Our system
forms a closed loop: \visiscore assesses quality; the agent decides whether to diagnose directly,
diagnose with caution, enhance-then-diagnose (\visienhance), or query for a retake; and a frozen
classifier produces the final diagnosis, its posterior confidence being the signal the gate routes on
(a standard bottleneck posterior suffices, \S\ref{sec:exp-dca}). Crucially, the driving quality scalar
need not come from \visiscore alone: any per-input estimate --- a learned scorer, a no-reference IQA
metric, or a known corruption severity --- keeps the framework general.

\paragraph{Contributions.}
To motivate quality-aware triage, we first chart a reliability~$\times$~recoverability decision
surface (per-dimension~$\times$~per-severity, AUC-led, with a recoverability axis derived from
enhancement) that quantifies where degradation breaks diagnosis and where enhancement recovers it.
Complementing rather than reusing prior per-dimension calibration views, it frames the three
contributions that follow.
\begin{itemize}
\item \textbf{C1 --- Diagnosis-preserving enhancement with a mutual-information lower bound.} We
introduce a feature-level diagnosis-preserving objective (DP-Loss) under which enhancement does not
erode diagnostic content within a moderate-degradation window. We derive a Pinsker-optimal
mutual-information lower bound, $I(\Zenh;Y)\geq I(\Zref;Y)-\beta\sqrt{\epsilon}$ under a KL budget
$\Ldp\leq\epsilon$ (Lemma~\ref{lem:dp}), together with a directional entropy result
(Proposition~\ref{prop:entropy}), and confirm them via the DP-Loss ablation
($\Delta\mathrm{AUC}_{\text{enh}}$ improves by $+0.0205$, $95\%$ CI $[+0.005,+0.035]$; McNemar
$p=2.3\times10^{-45}$).
\item \textbf{C2 --- A query-for-retake agent.} Our routing policy selects among four channels ---
direct diagnosis, cautioned diagnosis, enhance-then-diagnose, and query-for-retake --- of which, to
our knowledge, query-for-retake is the first to request \emph{re-acquisition} rather than merely
abstaining or deferring to a clinician. A local, per-band conditional risk guarantee
(Theorem~\ref{thm:agent-compact}) backs its licence to enhance: the benefit is positive only inside a
quality band $[\tauenh,\tauhigh]$, outside which the policy routes to a retake.
\item \textbf{C3 --- An honest reliability boundary.} We show, and do not hide, where enhancement turns
net-negative: severely degraded malignant cases, where reweighting the reconstruction loss fails to
close the salvage gap (a melanoma salvage of $5.2\%$ against a $31\%$ damage rate), and extreme
contrast loss at the per-dimension level. These boundaries are the hardest motivation for the
query-for-retake gate; correspondingly, under a fixed-threshold policy global triage does not yet beat
direct diagnosis (confidence-gated sensitivity $0.818$ versus $0.788$), which we report as motivation
for learned routing thresholds rather than a result to suppress.
\end{itemize}
